\documentclass[11pt]{article}
\usepackage[margin=1in]{geometry}
\usepackage{amsmath,amssymb}
\usepackage{booktabs}

\title{Methodology for Kernel Density Estimation via Diffusion}
\author{}
\date{}

\begin{document}
\maketitle

\section{Data}

Let \(X_1,\ldots,X_N\) be independent observations from an unknown density
\(f\). The goal is to estimate \(f\) nonparametrically while reducing boundary
bias and allowing adaptive smoothing through a diffusion model.

\section{Gaussian KDE}

The Gaussian kernel with smoothing parameter \(t>0\) is
\[
\phi(x,y;t)
=
\frac{1}{\sqrt{2\pi t}}
\exp\left(-\frac{(x-y)^2}{2t}\right).
\]
The ordinary Gaussian kernel density estimator is implemented as
\[
\widehat f(x;t)
=
\frac{1}{N}\sum_{i=1}^N \phi(x,X_i;t).
\]
No library KDE routine is used.

\section{Bandwidth Parameter}

The ordinary bandwidth is \(h=\sqrt{t}\). The project keeps three smoothing
parameters separate:
\[
\widehat t,\qquad \widehat t_2,\qquad t_\star.
\]
Here \(\widehat t\) builds the Gaussian pilot density, \(\widehat t_2\) is an
auxiliary bandwidth for estimating roughness, and \(t_\star\) is the final
diffusion time.

\section{Algorithm 1}

Algorithm 1 is the Improved Sheather--Jones plug-in selector from Botev,
Grotowski and Kroese. For derivative order \(j\), equation (26) estimates
\[
\|f^{(j)}\|^2
\approx
\frac{(-1)^j}{N^2}
\sum_{k=1}^N \sum_{m=1}^N
\phi^{(2j)}(X_k,X_m;2t_j).
\]
The recursive map \(\gamma_j\) is based on equation (29):
\[
\widehat t_j
=
\left[
\frac{1+2^{-(j+1/2)}}{3}
\frac{1\cdot 3\cdot 5\cdots(2j-1)}
{N\sqrt{\pi/2}\,\widehat{\|f^{(j+1)}\|^2}}
\right]^{2/(3+2j)}.
\]
The fixed-point iteration uses
\[
z_{n+1}
=
\xi\,\gamma^{[l]}(z_n),
\qquad
\xi
=
\left(\frac{6\sqrt 2-3}{7}\right)^{2/5},
\qquad l=5.
\]
The implementation evaluates the derivative functionals using the DCT
acceleration recommended by the paper, but the recursion itself is implemented
directly.

\section{Outputs \(\widehat t\) and \(\widehat t_2\)}

Algorithm 1 returns \(\widehat t\), the bandwidth parameter for the Gaussian
pilot KDE, and \(\widehat t_2\), the auxiliary bandwidth associated with
estimating \(\|f''\|^2\). The implementation also records whether the
fixed-point iteration converged and how many iterations were used.

\section{Pilot Density}

After Algorithm 1, the pilot density is
\[
p(x)
=
\frac{1}{N}\sum_{i=1}^N \phi(x,X_i;\widehat t).
\]
This density is only a pilot for the diffusion operator. It is not the final
estimator.

\section{Diffusion PDE}

The diffusion estimator solves
\[
\frac{\partial g}{\partial t}
=
\frac{1}{2}\frac{\partial}{\partial x}
\left[
a(x)\frac{\partial}{\partial x}
\left(\frac{g(x,t)}{p(x)}\right)
\right],
\]
with empirical initial condition \(g(x,0)=N^{-1}\sum_i\delta(x-X_i)\).
On a bounded interval, no-flux boundary conditions are used.

\section{Choice \(a(x)=p(x)\)}

The numerical experiments use \(a(x)=p(x)\), as requested. Therefore
\[
\sigma^2(x)=\frac{a(x)}{p(x)}=1,
\qquad
\sigma(x)=1.
\]
The finite-volume discretization uses cell-centered densities and internal
fluxes
\[
J_{i+1/2}
=
-\frac{1}{2}a_{i+1/2}
\frac{r_{i+1}-r_i}{\Delta x},
\qquad
r_i=\frac{g_i}{p_i}.
\]
The boundary fluxes are set to zero, preserving total probability mass.

\section{Estimating \(\|Lf\|^2\)}

The paper's plug-in idea is used:
\[
\|Lf\|^2
\approx
\|Lg(\cdot;\widehat t_2)\|^2.
\]
Since \(Lg=\partial g/\partial t\), the implementation estimates
\[
\frac{\partial g}{\partial t}(x;\widehat t_2)
\approx
\frac{
g(x;\widehat t_2+\varepsilon_t)
-
g(x;\widehat t_2)
}{\varepsilon_t}.
\]
Then numerical quadrature gives
\[
\widehat R
=
\int
\left[
\frac{\partial g}{\partial t}(x;\widehat t_2)
\right]^2 dx.
\]

\section{Calculating \(t_\star\)}

The asymptotic formula is
\[
t_\star
=
\left[
\frac{\mathbb E_f[\sigma^{-1}(X)]}
{2N\sqrt{\pi}\,\|Lf\|^2}
\right]^{2/5}.
\]
Because \(a(x)=p(x)\), \(\sigma(x)=1\), so
\[
t_\star
=
\left[
\frac{1}
{2N\sqrt{\pi}\,\widehat R}
\right]^{2/5}.
\]

\section{Final Diffusion Estimate}

After \(t_\star\) is computed, the solver returns to the empirical initial
condition and evolves the diffusion equation from \(t=0\) to \(t=t_\star\).
The final estimator is
\[
g(x;t_\star).
\]

\section{ISE}

When the true density is known, accuracy is measured by integrated squared
error:
\[
\mathrm{ISE}
=
\int
\left[
\widehat f(x)-f(x)
\right]^2 dx.
\]
The code evaluates this integral numerically on the plotting or diffusion
grid.

\section{Figure 1 Reproduction}

The Figure 1 reproduction uses
\[
f(x)=4(1-x)^3,\qquad 0\le x\le 1,
\]
which is the \(\mathrm{Beta}(1,4)\) density, with \(N=1000\). The bandwidth is
fixed from the paper:
\[
\sqrt t=0.05248.
\]
This experiment does not run Algorithm 1. It compares the true density, the
manual Gaussian KDE, and the reflected boundary estimator
\[
\kappa(x,X_i;t)
=
\sum_{k=-\infty}^{\infty}
\left[
\phi(x,2k+X_i;t)+\phi(x,2k-X_i;t)
\right].
\]
The infinite reflection sum is truncated numerically when additional terms are
negligible.

\section{Sample-Size Variation}

The variation experiment uses the same \(\mathrm{Beta}(1,4)\) target density and
sample sizes
\[
N\in\{50,100,250,500,1000,5000\}.
\]
For each \(N\), independent samples are generated, Algorithm 1 constructs the
pilot bandwidths, Algorithm 2 computes the diffusion KDE, and ISE is recorded.
The script saves both summary and repetition-level CSV files and plots mean
ISE against sample size.

\end{document}
